% !TEX root = main.tex
\documentclass[main.tex]{subfiles}
\begin{document}

\section{Constrained Power Diagrams}
\label{sec:contributions}

In this section, we present the construction of power diagrams whose cells have prescribed centroids $\smash{\tcentroid_i \!\in\! \R^d}$, collected as $\smash{\tcentroid \!=\! \{\tcentroid_i\}_{i=1}^n}$. %, chosen independently of the sites.
Similar to the volume-constrained power diagrams in \eqref{eq:volume-dual}, we show that these centroid-constrained power diagrams are generated by maximizing an unconstrained concave objective, which yields a unique partition whenever one exists, both for centroid constraints alone (\S\ref{sec:centroid-constraints}) and jointly with volume constraints (\S\ref{sec:joint-constraints}).
We further propose a softening of the centroid constraints to handle targets that no power diagram realizes (\S\ref{sec:regularized-joint}), and close with remarks (\S\ref{sec:remarks}).

\subsection{Centroid Constraints}
\label{sec:centroid-constraints}

We first consider the construction of a partition $\cell$ associated with fixed weighted sites $(\qq, \weight)$ whose cells have prescribed centroids $\tcentroid$.
Reusing the objective from \eqref{eq:power-diagram}, we formulate this problem as
\begin{subequations}
\label{eq:centroid-constrained}
\begin{alignat}{2}
  & \min_{\cell} \quad && \sum\nolimits_{i} \int_{\cell_i} \bigl( \pdist{\xx}{i} \bigr) \, d\xx \label{eq:centroid-constrained-obj} \\
  & \text{s.t.} \quad && \int_{\cell_i} (\xx - \tcentroid_i) \, d\xx = 0 \ \ \forall i. \label{eq:centroid-constrained-cons}
\end{alignat}
\end{subequations}
% 
Introducing a vector multiplier $\shift_i \!\in\! \R^d$, which we call the \emph{shift}, for each centroid constraint \eqref{eq:centroid-constrained-cons} yields the Lagrangian
\begin{equation}
\lag(\cell, \shift)
\!=\! 
\sum\nolimits_{i} \int_{\cell_i}\! \bigl( \pdist{\xx}{i} \bigr) d\xx - 2 \shift_i\tr\! \int_{\cell_i}\! (\xx - \tcentroid_i) \, d\xx,
\label{eq:centroid-lagrangian}
\end{equation}
% 
where we indicate the collected shifts as $\smash{\shift\!=\!\{\shift_i\}_{i=1}^n}$ and scale each $\shift_i$ by a factor of two to make our derivatives more concise.
% 
For fixed shifts, the combined integrand in \eqref{eq:centroid-lagrangian} can be rearranged as
\begin{equation}
  \pdist{\xx}{i} - 2\,\shift_i\tr (\xx - \tcentroid_i)
    = \norm{\xx - \dual_i}^2 - \combweight_i ,
  \label{eq:complete-square}
\end{equation}
based on the \emph{shifted} weighted sites $(\dual_i, \combweight_i)$ of the form
\begin{subequations}
  \label{eq:shifted-site}
  \begin{empheq}[left={\empheqlbrace\,}]{align}
    \dual_i &= \qq_i + \shift_i , \label{eq:shifted-site-pos} 
    \\[2mm]
    \combweight_i &= \weight_i + \norm{\shift_i}^2 + 2\,\shift_i\tr (\qq_i - \tcentroid_i). \label{eq:shifted-site-weight}
  \end{empheq}
\end{subequations}
% 
We thus conclude that, for fixed shifts, minimizing \eqref{eq:centroid-lagrangian} over partitions yields the power diagram of the shifted weighted sites \eqref{eq:shifted-site}.
% 
We call the shifted sites $\dual$ the \emph{generators} of this power diagram, to distinguish them from the sites $\qq$ used in the objective \eqref{eq:centroid-constrained-obj}.
Importantly, the properties of power diagrams reviewed in \S\ref{sec:background}, such as the homotheties $\homotheties$ and the normals of $\achievable$, apply to the generators $\dual$ rather than to the sites $\qq$.
% 
As in \S\ref{sec:background}, we restrict the shifts to configurations with pairwise distinct generators $\dual$ and non-empty cells, which guarantees that this minimizer is a unique power diagram. % on which $\centroiddual$ is differentiable.

The optimal value of \eqref{eq:centroid-lagrangian} defines the \emph{dual objective}
\begin{equation}
  \centroiddual(\shift)
  % =
  % \centroiddual(\shift; \qq, \weight, \tcentroid)
  = 
  \min_{\cell} \lag(\cell, \shift)
  = 
  \energy\bigl( \dual(\shift), \combweight(\shift) \bigr).
  \label{eq:centroid-dual}
\end{equation}
The constrained optimization in \eqref{eq:centroid-constrained} thus becomes the {unconstrained} maximization of $\centroiddual$ over the space of power diagrams, parametrized by the shifts $\shift$.
% , with the weighted sites $(\qq, \weight)$ and the targets $\tcentroid$ as fixed inputs.
% 
As detailed in Appendix~\ref{sec:appendix-centroid}, the gradient of $\centroiddual$ recovers the centroid constraint
\begin{equation}
  \partial_{\shift_i} \centroiddual
    = 2 \int_{\cell_i} (\tcentroid_i - \xx) \, d\xx
    = 2 \, \cellvol_i \, (\tcentroid_i - \centroid_i) ,
  \label{eq:centroid-dual-gradient}
\end{equation}
so a stationary shift is exactly one whose cells realize the prescribed centroids.
The Hessian of $\centroiddual$ is a (negated) facet-based vector Laplacian coupling neighboring shifts through the facets $\face_{ij}$ of the power diagram with $d \!\times\! d$ blocks
\begin{subequations}
  \label{eq:centroid-hessian}
  \begin{empheq}[left={\empheqlbrace\,}]{align}
    \partial^2_{\shift_i \shift_j}
    \centroiddual &= \frac{2}{\norm{\dual_i - \dual_j}} \int_{\face_{ij}} (\xx - \tcentroid_i)(\xx - \tcentroid_j)\tr \, dA,
    \label{eq:centroid-hessian-off}
    \\[2mm]
    \partial^2_{\shift_i \shift_i}
    \centroiddual &= -\!\sum_{j \sim i} \frac{2}{\norm{\dual_i - \dual_j}} \int_{\face_{ij}} (\xx - \tcentroid_i)(\xx - \tcentroid_i)\tr \, dA.
    \label{eq:centroid-hessian-diag}
  \end{empheq}
\end{subequations}
In Appendix~\ref{sec:appendix-concavity}, we show that \eqref{eq:centroid-hessian} is \emph{negative semi-definite} and, hence, $\centroiddual$ is concave with respect to shifts.
Therefore, the centroid-constrained problem is an \emph{unconstrained concave maximization}.
% 
As discussed in Appendix~\ref{sec:appendix-nullspace}, the Hessian of $\centroiddual$ may also have a nontrivial null-space, whose modes move the generators while keeping the weights fixed, but only for special centroid configurations.
Either way, this concerns the uniqueness of the \emph{shifts}, while the resulting power diagram is unique whenever it exists.
% any partition solving \eqref{eq:centroid-constrained} minimizes the Lagrangian \eqref{eq:centroid-lagrangian} at every maximizing shift, and this minimizer is unique.
% We discuss the existence of maximizers, and of maximizers with distinct shifted sites, in \S\ref{sec:remarks}.

\subsection{Volume + Centroid Constraints}
\label{sec:joint-constraints}

Next, we consider the partition that enforces both prescribed volumes $\smash{\pvol}$ \emph{and} prescribed centroids $\smash{\tcentroid}$ simultaneously.
Repeating the steps of \S\ref{sec:centroid-constraints}, we introduce a shift $\shift_i \!\in\! \R^d$ for each centroid constraint and a weight $\weight_i \!\in\! \R$ for each volume constraint, which together yield the Lagrangian
\begin{equation}
  \lag(\cell, \shift, \weight)
  \!=\! 
  \sum_{i} 
  \!
  \int_{\cell_i} \!\!\!\! \norm{\xx -\! \qq_i}^2 \!d\xx 
  - 
  2\shift_i\tr 
  \!\!
  \int_{\cell_i} \!\!\!\! (\xx - \tcentroid_i) d\xx 
  - 
  \weight_i \bigl( \cellvol_i - \pvol_i \bigr).
  \label{eq:joint-lagrangian}
\end{equation}
For fixed multipliers, the partition minimizer corresponds once again to a power diagram with shifted weighted sites $(\dual, \combweight)$ given in \eqref{eq:shifted-site}, under the same restriction to distinct shifted sites with non-empty cells as in \S\ref{sec:centroid-constraints}.
Its optimal value defines the \emph{dual objective}
\begin{equation}
  \jointdual(\shift, \weight)
  % \jointdual(\shift, \weight; \qq, \tcentroid)
    \!=\! 
    \min_{\cell} \lag(\cell, \shift, \weight)
    \!=\!
    \centroiddual(\shift, \weight) \!+\! \sum\nolimits_{i}\! \weight_i \pvol_i ,
  \label{eq:joint-dual}
\end{equation}
where $\smash{\centroiddual(\shift, \weight) = \energy\bigl( \dual(\shift,\weight), \combweight(\shift,\weight) \bigr)}$ is the centroid dual \eqref{eq:centroid-dual}, now with the weights promoted from fixed inputs to unknowns.
The constrained problem thus becomes the {unconstrained} maximization of $\jointdual$ over the space of power diagrams, parametrized by both the shifts $\shift$ and weights $\weight$.
% now with the sites $\qq$ and the prescribed centroids $\tcentroid$ and volumes $\pvol$ as fixed inputs.

Following Appendix~\ref{sec:appendix-joint}, the gradient of $\jointdual$ stacks the centroid and volume residuals
\begin{subequations}
  \label{eq:joint-gradient}
  \begin{empheq}[left={\empheqlbrace\,}]{align}
    \partial_{\shift_i} \jointdual &= 2\,\cellvol_i\,(\tcentroid_i - \centroid_i) , \label{eq:joint-gradient-shift} 
    \\[2mm]
    \partial_{\weight_i} \jointdual &= \pvol_i - \cellvol_i , \label{eq:joint-gradient-weight}
  \end{empheq}
\end{subequations}
so a stationary configuration is one whose cells realize both the prescribed centroids and volumes.
% 
The Hessian of $\jointdual$ is again a (negated) facet-based Laplacian, now assembled from the two multiplier types per site.
Its {diagonal} blocks are exactly the two Hessians we have already seen: the shift--shift block reproduces \eqref{eq:centroid-hessian} from \S\ref{sec:centroid-constraints}, and the weight--weight block reproduces the volume Hessian \eqref{eq:newton-step} from \S\ref{sec:background} for the power diagram of the shifted sites, \ie with $\smash{\norm{\qq_i - \qq_j}}$ replaced with $\smash{\norm{\dual_i - \dual_j}}$.
The only new term is the mixed shift--weight block, which takes the form of:
\begin{subequations}
  \label{eq:joint-hessian}
  \begin{empheq}[left={\empheqlbrace\,}]{align}
    \partial^2_{\shift_i \weight_j} \jointdual
      &= \frac{1}{\norm{\dual_i - \dual_j}} \int_{\face_{ij}} (\xx - \tcentroid_i) \, dA,
    \label{eq:joint-hessian-off}
    \\[2mm]
    \partial^2_{\shift_i \weight_i} \jointdual
      &= -\!\sum\nolimits_{j \sim i} \partial^2_{\shift_i \weight_j} \jointdual .
    \label{eq:joint-hessian-diag}
  \end{empheq}
\end{subequations}
Since the weight gradient in \eqref{eq:joint-gradient-weight} is the volume residual, the mixed block in \eqref{eq:joint-hessian} is the negated gradient of the cell volumes with respect to the shifts, \ie $\smash{\partial^2_{\shift_i \weight_j} \jointdual = -\partial \cellvol_j / \partial \shift_i}$.
This mixed term can thus be interpreted as a discrete divergence operator as in \citet{deGoes2015PowerParticles}, which quantifies the infinitesimal volume change of the cells as the shifted sites $\dual$ are displaced by infinitesimal changes of the shifts.
Accordingly, its diagonal \eqref{eq:joint-hessian-diag} states that these volume changes sum to zero, since the cells always tile $\domain$.

Similar to $\centroiddual$, the Hessian of $\jointdual$ is \emph{negative semi-definite}, so $\jointdual$ is jointly concave in the shifts and weights, and the joint problem is again an \emph{unconstrained concave maximization} (Appendix~\ref{sec:appendix-concavity}).
Unlike $\centroiddual$, the Hessian of $\jointdual$ always has a nontrivial null-space, since the free weights let each multiplier $\smash{(\shift_i, \weight_i)}$ move its lifted site freely in $\smash{\R^{d+1}}$.
In particular, this null-space contains every infinitesimal homothety of the shifted weighted sites, \ie the first-order variation of \eqref{eq:lifted-homothety} about the identity.
In terms of the multipliers, these infinitesimal homotheties are the perturbations
\begin{equation}
  \delta\shift_i = \kappa\,\dual_i + \translation , \qquad \delta\weight_i = \kappa\,\gamma_i + 2\,\translation\tr \tcentroid_i - c ,
  \label{eq:joint-nullspace}
\end{equation}
for any $\smash{(\kappa, \translation, c) \in \R \times \R^d \times \R}$, where
\begin{equation}
  \gamma_i = \norm{\tcentroid_i}^2 - \bigl( \norm{\tcentroid_i - \qq_i}^2 - \weight_i \bigr)
  \label{eq:pivot}
\end{equation}
is the height at the target $\tcentroid_i$ of the hyperplane lifted from the weighted site $(\qq_i, \weight_i)$, which does not depend on the shift $\shift_i$.
These perturbations reproduce  the infinitesimal version of the redundancy of the triplets $\smash{(\qq_i, \shift_i, \weight_i)}$ derived by \citet[Lemma~1]{Memari2012PrimalDual}.
% 
As detailed in Appendix~\ref{sec:appendix-nullspace}, these $d \!+\! 2$ modes cover the entire Hessian null-space of $\jointdual$, except when the dual graph has cut cells or the generators are not in general position, which can then add localized homotheties.
% at the cut cells and at cycles of more than three dual edges, respectively.
% 
Note also that the perturbations whose weight updates vanish are the null-space modes of $\centroiddual$ and, for $\kappa \!=\! 0$ and $\translation \!=\! 0$, they reduce to the constant spanning the Hessian null-space of $\dualenergy$.
% 
Since all null-space modes leave the cells unchanged, the maximizing multipliers, when they exist, generate a \emph{single} power diagram.

\subsection{Softened Centroid Constraints}
\label{sec:regularized-joint}

Although the jointly constrained problem in \S\ref{sec:joint-constraints} has a unique solution whenever one exists, not every choice of volume and centroid targets is realized by a power diagram.
% since realizable targets must be exposed points of the achievable centroids $\achievable$ 
% 
However, the two constraint types play asymmetric roles in this limitation.
Any positive volumes verifying \eqref{eq:partition-feasibility-vol} are realized by a power diagram for any distinct sites \cite{Aurenhammer1998Minkowski}, so infeasibility arises only once the centroids are prescribed as well.
% , as previously characterized by the set $\achievable$ of achievable centroids in \S\ref{sec:background}.
This asymmetry motivates us to consider partitions of $\domain$ in which the volume constraints remain exact and only the centroid constraints are {softened}.

The most direct way to soften the centroid constraints is to penalize the least-squares centroid residual, as done by \citet{Bourne2025Inverting}, who minimized this penalty over the sites while enforcing the volumes through the weights.
However, as acknowledged by the authors, this objective is non-convex in the sites and its optimization is challenging, with local solutions found by a convex surrogate only under implicit assumptions on the target centroids.
In contrast, we show next that the formulation of \S\ref{sec:joint-constraints} admits a small modification that incorporates the softening while keeping the problem an unconstrained concave maximization.

We start from the problem underlying \S\ref{sec:joint-constraints}, which minimizes the transport cost from sites $\qq$ to cells $\cell$ subject to centroid and volume constraints.
We relax each centroid constraint by a slack variable $\smash{\slack_i \!\in\! \R^d}$ that absorbs its residual, and penalize the squared norm of these slacks in the objective, yielding a modified constrained optimization of the form
\begin{subequations}
  \label{eq:regularized-primal}
  \begin{alignat}{3}
    & \min_{\cell, \slack} \quad & & \sum\nolimits_{i} \int_{\cell_i} \norm{\xx - \qq_i}^2 \, d\xx + \alpha \sum\nolimits_i \pvol_i \norm{\slack_i}^2 \label{eq:regularized-primal-obj} \\
    & \text{s.t.} \quad & \constraintbrace{16pt}{8pt} & \int_{\cell_i} (\xx - \tcentroid_i) \, d\xx = \pvol_i \, \slack_i \ \ \forall i , \label{eq:regularized-primal-cons} \\
    & & & \cellvol_i = \pvol_i \ \ \forall i , \label{eq:regularized-primal-vol}
  \end{alignat}
\end{subequations}
where $\alpha \!>\! 0$ acts as the stiffness of the penalty.
% As special cases, note that $\alpha \!=\! 0$ leaves the slacks free thus reducing to the volume-constrained problem \eqref{eq:volume-constrained}, while $\alpha \!\to\! \infty$ forces $\slack \!=\! 0$ and recovers the joint problem of \S\ref{sec:joint-constraints}.
Introducing the same multipliers as in \S\ref{sec:joint-constraints}, the optimal slacks take the closed form $\smash{\slack_i = -\shift_i / \alpha}$, and substituting them back into the Lagrangian leads to the \emph{softened dual objective}
\begin{equation}
  % \jointdual_{\alpha}(\shift, \weight; \qq, \tcentroid) 
  \jointdual_{\alpha}(\shift, \weight) 
  = 
  % \jointdual(\shift, \weight; \qq, \tcentroid) 
  \jointdual(\shift, \weight) 
  - 
  \frac{1}{\alpha} \sum\nolimits_i \pvol_i \norm{\shift_i}^2 ,
  \label{eq:regularized-dual}
\end{equation}
which only adds a concave quadratic term on the shifts to the joint dual \eqref{eq:joint-dual}.
The softened problem is thus again the \emph{unconstrained maximization of a concave function} over the shifts and weights.

Compared to \S\ref{sec:joint-constraints}, the penalty only modifies the shift derivatives of $\jointdual$.
The weight gradient \eqref{eq:joint-gradient-weight} is unchanged, so the volumes remain exact at a maximizer, while the shift gradient \eqref{eq:joint-gradient-shift} becomes
\begin{equation}
  \partial_{\shift_i} \jointdual_{\alpha} = 2\,\cellvol_i\,(\tcentroid_i - \centroid_i) - \frac{2}{\alpha}\,\pvol_i\,\shift_i .
  \label{eq:regularized-gradient-shift}
\end{equation}
Likewise, the Hessian of \eqref{eq:regularized-dual} keeps the weight--weight and mixed blocks from \eqref{eq:joint-hessian}, while each diagonal shift--shift block is shifted by $\smash{-\tfrac{2}{\alpha}\,\pvol_i\,\Id}$.
% Since the Hessian of $\jointdual$ is negative semi-definite, 
This extra term makes the Hessian of $\jointdual_\alpha$ strictly negative along any perturbation with a nonzero shift component, which removes every null-space mode of $\jointdual$ that moves the shifts, \ie the homotheties \eqref{eq:joint-nullspace} as well as any localized homothety (Appendix~\ref{sec:appendix-nullspace}).
The remaining perturbations change only the weights, along which the Hessian reduces to the volume Laplacian \eqref{eq:newton-step} of the shifted sites, so the null-space consists only of the additive constant on the weights.
Consequently, $\jointdual_{\alpha}$ is concave for any $\alpha \!>\! 0$ and admits a unique maximizer for {any} target centroids.

While the maximizer of $\jointdual_\alpha$ always exists, the resulting partition is a power diagram only if its generators $\smash{\dual}$ remain pairwise distinct.
% We derive next a bound on $\alpha$ that ensures distinct generators.
To this end, we derive in Appendix~\ref{sec:appendix-softened-bound} an upper bound on $\alpha$ that guarantees they remain distinct, which is written as
\begin{equation}
  \alpha < \min_{i \neq j} \frac{\norm{\qq_i - \qq_j}}{\sqrt{\resid\,\bigl(1/\pvol_i + 1/\pvol_j\bigr)}} ,
  \label{eq:regularized-distinct}
\end{equation}
where $\smash{\resid = \sum_i \pvol_i \norm{\centroid^0_i - \tcentroid_i}^2}$ is the residual between the prescribed centroids $\tcentroid$ and the centroids $\smash{\centroid^0}$ obtained by the power diagram that maximizes $\dualenergy$ in \eqref{eq:volume-dual} with sites $\qq$ and prescribed volumes $\pvol$.
Note that \eqref{eq:regularized-distinct} is conservative, since it bounds the residual of each cell by the residual $\resid$ of all cells.
% 
In our experiments, however, the volume constraints rarely let the residual of one cell move distant cells, so we also propose a practical, but not guaranteed, bound estimate:
\begin{equation}
  \alpha < \min_{i \neq j} \frac{\norm{\qq_i - \qq_j}}{\norm{\centroid^0_i - \tcentroid_i} + \norm{\centroid^0_j - \tcentroid_j}} .
  \label{eq:regularized-practical}
\end{equation}

Moreover, observe that, at the maximizer, \eqref{eq:regularized-primal-cons} forces $\smash{\slack_i = \centroid_i \!-\! \tcentroid_i}$, and the optimal shifts follow from \eqref{eq:regularized-gradient-shift} as $\smash{\shift_i = \alpha\,(\tcentroid_i \!-\! \centroid_i)}$.
The penalty in \eqref{eq:regularized-primal-obj} is thus the volume-weighted least-squares centroid residual $\smash{\alpha \sum_i \pvol_i \norm{\centroid_i - \tcentroid_i}^2}$, which resembles the residual of \citet{Bourne2025Inverting}, except that they weighted each cell by $\smash{\pvol_i^2}$ instead of $\pvol_i$.
While they minimized their residual through a non-convex search over the sites, we optimize ours through a concave maximization.
% Our softened problem penalizes this same residual combined with the first term of \eqref{eq:regularized-primal-obj} through a concave maximization.
% At the solution, we can further rearrange \eqref{eq:regularized-primal-obj}, up to a constant, as
Expanding each integral about its cell centroid, with exact volumes, the optimal value of the objective \eqref{eq:regularized-primal-obj} further decomposes as
\begin{equation}
  \sum\nolimits_i \int_{\cell_i} \norm{\xx \!-\! \centroid_i}^2 d\xx
  + \sum\nolimits_i \pvol_i \norm{\centroid_i \!-\! \qq_i}^2
  + \alpha \sum\nolimits_i \pvol_i \norm{\centroid_i \!-\! \tcentroid_i}^2 ,
  \label{eq:regularized-decomposition}
\end{equation}
\ie the quantization error of the cells plus two least-squares terms that pull the centroids towards the sites and towards the targets, respectively.
Expanding the quantization error and dropping the terms independent of the partition, \eqref{eq:regularized-decomposition} rearranges as
\begin{equation}
  \alpha \sum\nolimits_i \pvol_i \norm{\centroid_i - \bigl( \tcentroid_i + \tfrac{1}{\alpha}\,\qq_i \bigr)}^2.
  \label{eq:regularized-projection}
\end{equation}
Therefore, maximizing $\jointdual_\alpha$ returns the partition whose centroids $\centroid$ are the closest to the \emph{pushed targets} $\smash{\tcentroid + \qq/\alpha}$ within $\achievable$ in the volume-weighted least-squares sense.
% The generators of the resulting power diagram also follow from \eqref{eq:regularized-gradient-shift} as
% \begin{equation}
%   \dual_i = \qq_i + \shift_i = \alpha \Bigl( \tcentroid_i + \tfrac{1}{\alpha}\,\qq_i - \centroid_i \Bigr),
%   \label{eq:regularized-generators}
% \end{equation}
% indicating the difference between the pushed targets and their projection, which is an outward normal of $\achievable$ at the centroids $\centroid$ (\S\ref{sec:background}).

% Finally, we note that maximizing $\jointdual_\alpha$ corresponds to a proximal step on $\jointdual$ over the generators $\dual$, centered at the sites $\qq$.
% We can thus recover the maximizer of $\jointdual$ by iterating the maximization of $\jointdual_\alpha$, as in the proximal point algorithm \cite{Rockafellar1976Proximal}.
% Starting from sites at the targets $(\smash{\qq^0 \!=\! \tcentroid})$, we move the sites to the generators after each maximization $(\smash{\qq^{k+1} \!=\! \dual^k})$, and warm-start the next one at zero shifts and the previous shifted weights $(\smash{\weight^{k+1} \!=\! \combweight^k})$, which reproduces the previous power diagram.
% For realizable targets, the sites converge to a positive dilation of the generators realizing $\tcentroid$.
% In particular, a single maximization of $\jointdual_\alpha$ suffices when the sites already generate the power diagram realizing $\tcentroid$ up to a homothety \eqref{eq:lifted-homothety}, since the pushed targets then project onto $\tcentroid$ and \eqref{eq:regularized-generators} returns $\dual = \qq$.
% To keep every iterate a power diagram, we also decrease $\alpha$ at each iteration, if needed, so that it satisfies \eqref{eq:regularized-distinct} for the current sites. % (Lemma~\ref{lem:proximal-halving}).
% We summarize this scheme in Algorithm~\ref{alg:proximal}.

\subsection{Remarks}
\label{sec:remarks}

We close this section with a few remarks on our formulations.
We first discuss the feasibility of the prescribed values and the existence of solutions, then relate our results to \citet{Bourne2025Inverting, Bourne2026ConvexOrder}, and finally discuss the special case of reducible power diagrams.

\paragraph{Feasibility and existence.}
The formulations of \S\ref{sec:centroid-constraints} and \S\ref{sec:joint-constraints} require target values compatible with the conservation identities \eqref{eq:partition-feasibility}.
In the joint case, these identities can be checked directly on the data, while the centroid-only case requires finding positive volumes that satisfy \eqref{eq:partition-feasibility-centroid}, \eg with the construction of Appendix~\ref{sec:appendix-maxent}.
These tests are only necessary, not sufficient.
The targets must also be realizable, otherwise both duals may approach their supremum as generators merge. %, without attaining it.
In contrast, the softened formulation needs no test on the centroids, since its slacks absorb any mismatch with the targets.
Provided the volumes verify \eqref{eq:partition-feasibility-vol}, maximizing $\jointdual_{\alpha}$ always returns the partition whose centroids are the projection of the pushed targets onto $\achievable$, which is a power diagram under \eqref{eq:regularized-distinct}.

\paragraph{Link to \citet{Bourne2025Inverting}.}
Our joint problem in \S\ref{sec:joint-constraints} is closely related to the inverse problem considered by \citet{Bourne2025Inverting}, which recovers a power diagram from the volumes and centroids of its cells.
In their setting, the unknowns are the sites of the diagram, so the generators of the sought diagram and the sites used in the objective are the same.
In our formulation, the sites $\qq$ and targets $\tcentroid$ are prescribed independently, hence the generators of our solution are the shifted sites $\dual$ of \eqref{eq:shifted-site}.
% 
The two formulations also handle the volume and centroid constraints differently.
The objective gradient proposed by \citet{Bourne2025Inverting} equals the centroid residual only on the submanifold of weights that enforce the target volumes, so their characterization of solutions holds only once the volume constraints are satisfied, requiring every centroid update to be preceded by a weight solve.
We instead impose both constraints on an equal footing through Lagrange multipliers, the weights for the volumes and the shifts for the centroids.
For every value of the multipliers, the volume and centroid residuals are exactly the gradients of $\jointdual$ in \eqref{eq:joint-gradient}, so our characterization of solutions, as stationary points of a single concave function, holds without first satisfying either constraint.
In particular, the centroid residual can be reduced with the weights held fixed, which is precisely the centroid-only problem of \S\ref{sec:centroid-constraints}.
% 
Yet, the two formulations share the same solutions.
At weights that enforce the volumes, $\jointdual$ reduces to their objective, up to a factor of two and an additive constant. % (Appendix~\ref{sec:appendix-regularized}).
Therefore, our results generalize \citet{Bourne2025Inverting} by recovering their inverse problem when the volumes are enforced, while also admitting arbitrary sites, the centroid-only problem of \S\ref{sec:centroid-constraints}, and the softened extension of \S\ref{sec:regularized-joint}.

\paragraph{Link to \citet{Bourne2026ConvexOrder}.}
The progression from \S\ref{sec:joint-constraints} to \S\ref{sec:regularized-joint} also relates to the subsequent work of \citet{Bourne2026ConvexOrder}.
Their method dilates the target centroids $\tcentroid$ by a factor $s \!>\! 1$, projects these dilated targets $s\,\tcentroid$ onto $\achievable$, and takes the residual $s\,\tcentroid \!-\! \centroid$ to their projection $\centroid$ as the generators of a power diagram with centroids $\centroid$.
% This method also applies to unrealizable targets, yielding power diagrams whose centroids approximate them, when the resulting generators are distinct.
% 
Instead, our softening of the centroid constraints in \S\ref{sec:regularized-joint} projects the pushed targets $\smash{\tcentroid \!+\! \qq / \alpha}$ onto $\achievable$. % as in \eqref{eq:regularized-projection}.
In the special case of $\tcentroid_i \!=\! \qq_i$, maximizing $\jointdual_\alpha$ is therefore the same problem as the projection-dilation scheme of \citet{Bourne2026ConvexOrder} with dilation factor $s \!=\! 1 \!+\! 1/\alpha$, and returns the same power diagram.
While their method is solved by first-order iterations that each require a full volume solve, our approach offers a second-order Newton's method that updates both the shifts and the weights simultaneously.
% Moreover, since the sites of $\jointdual_\alpha$ are free, we can employ proximal iterations, as described in \S\ref{sec:regularized-joint}, to recover realizable target centroids exactly, whereas \citet{Bourne2026ConvexOrder} only approaches realizable targets in the limit as $s \!\to\! 1$.

% \paragraph{Softening the volumes.}
% For completeness, the volume constraints can be softened in the same way, by a scalar slack per cell with stiffness $\beta \!>\! 0$ and reference weights $\bweight$, as detailed in Appendix~\ref{sec:softened-volume}.
% The resulting dual subtracts the concave quadratic $\smash{\tfrac{1}{4\beta} \sum_i \pvol_i\,(\weight_i - \bweight_i)^2}$ from $\jointdual_\alpha$, so it is strongly concave in all multipliers and has a unique maximizer for any targets.
% Its optimal weights follow an equation of state, and thus act as a pressure with $\bweight$ as reference, and it recovers every formulation of this section as a limit of the stiffnesses $(\alpha, \beta)$.

\paragraph{Reducible diagrams.}
\citet[\S3.4]{Bourne2026ConvexOrder} call a power diagram \emph{reducible} when its cells can be grouped into the cells of a coarser power diagram, and characterize the irreducible diagrams by where their centroids lie among the exposed points of the set $\achievable$ of achievable centroids.
However, their criterion cannot decide from the cells of a given diagram whether it is reducible.
We provide such a criterion by relating reducibility to the Hessian null-space of $\jointdual$.
As detailed in Appendix~\ref{sec:appendix-nullspace-reducible}, we show that, for a convex domain $\domain$, a power diagram is reducible if and only if the Hessian of $\jointdual$ has null-space modes beyond the global homotheties of \eqref{eq:joint-nullspace}, and we characterize these modes from the graph dual to the power diagram.
For generators in general position, a power diagram is reducible exactly when it has a \emph{cut cell}, \ie a cell whose removal disconnects the dual graph (Figure~\ref{fig:appendix-reducible} top).
Without general position, reducibility may also arise from non-triangular cycles of the dual graph (Figure~\ref{fig:appendix-reducible} bottom).
These additional modes are \emph{localized homotheties}, which apply a homothety to the lifted generators of only a part of the power diagram, and thus move the generators without changing the cells.
Therefore, the maximizers of $\jointdual$ form a convex set that extends beyond the global homotheties.
% , and that is bounded by configurations in which some generators coincide.
% In Figure~\ref{fig:appendix-reducible}, we provide examples for both cases, namely a diagram reducible through a cut cell (top row), and a symmetric diagram with vertices shared by four or eight cells (bottom row).

\end{document}
